\begin{corollary}[Sub-arrays and super-sequences]\label{lem:subarr}
Let $r$ be a quasi-order on $Q$.  The following are equivalent:
\begin{enumerate}
\item $Q$ is a better-quasi-order, i.e. $\operatorname{Bqo}(r)$;
\item every locally constant multi-sequence in $Q$ has a restriction to an
      infinite subbase which is perfect;
\item every super-sequence $f:F\to Q$, for a front $F$ on an arbitrary
      infinite base, has a restriction to a sub-front on an infinite
      subbase which is perfect.
\end{enumerate}
\end{corollary}
\begin{proof}
Assume first that $Q$ is bqo.  Let $h$ be locally constant.  Apply the
perfect-or-bad shift dichotomy \noderef{ShiftDichotomy}.  If its restricted
sequence were bad, then it would contradict the defining condition
\noderef{Bqo}, because local constancy is preserved by
\noderef{RestrictionLocallyConstant}.  Hence the perfect alternative holds,
which proves clause (2).

Now let $f:F\to Q$ be a super-sequence on an arbitrary infinite base $X$.
Its base extension is supplied by \noderef{SuperSequenceExtension}; it is
locally constant and agrees with $f$ on every selected front member.  The
base restriction lemma \noderef{BaseBqoRestriction} gives an increasing
restriction of this extension which is perfect.  Applying
\noderef{PerfectBaseToSuper}, and using \noderef{FrontRestriction} for the
resulting infinite subbase, yields a perfect restriction of $f$.  This proves
clause (3).

Conversely, assume clauses (2) and (3).  To prove $\operatorname{Bqo}(r)$,
let $h$ be locally constant and suppose for contradiction that it is bad.
The minimal deciding prefixes form a front by \noderef{DecidingFront}, and
\noderef{DecidingValue} supplies the induced value map
$h^{\dcl}:F^h\to Q$.  The badness transport
\noderef{BadMultiToSuper} makes this super-sequence bad.  Clause (3) gives
a perfect restriction to a sub-front; \noderef{BadSuperRestrict} keeps the
restriction bad.  This contradicts \noderef{NoBadPerfectSuper}.  Therefore
no locally constant multi-sequence is bad, which is exactly
\noderef{Bqo}.  All three clauses are consequently equivalent.
\end{proof}
